\centerline{\Huge \bf  Олимпиада I}
\centerline{\Huge \bf 7 класс}

\begin{flushright}
    \scriptsize{Заключительный отбор в Сириус 2021г.\\ Кто загуглит - тому бан!}
\end{flushright}

\problem{1} 
Из пункта $A$ в пункт $B$ выезжает Валя на велосипеде, а одновременно из пункта $B$ в пункт $A$ пешком выходит Галя. Через час, когда они встретились, Валя отдает Гале велосипед, и продол- жает движение в пункт $B$ пешком, а Галя устремляется в $A$ на велосипеде. Известно, что Валя пешком передвигается вдвое быстрее Гали, а на велосипеде Галя ездит в $1,5$ раза быстрее Вали. Валя пришла в пункт $B$ в $11:00$. В какое время Галя прибудет в пункт $A$?

\problem{2}
 В остроугольном треугольнике $ABC$ точки $D$ и $E$ $–$ основания высот, проведенных из вершин B и C соответственно. $\angle CBD = \angle ABD$ и $\angle ACE = \angle BCE$. Найдите углы треугольника $ABC$.

\problem{3}
 Если написать возраст Александра, а затем возраст Николая на сегодняшний день, то получится четырёхзначное число, которое является квадратом натурального числа. Если сделать то же самое через $11$ лет, то снова будет точный четырёхзначный квадрат. Найдите текущий возраст Александра и Николая.

\problem{4}
В выпуклом четырехугольнике $ABCD$ равны углы$ A, B, C$. Кроме того равны стороны $AD$ и $CD$. На стороне AB выбрана такая точка $E$, что $BE = AD.$ Докажите, что $CE$ $-$ биссектриса угла $BCD$.

\problem{5}
Незнайка очень обеспокоен своей средней оценкой по математике. Она оказалась меньше $3$. (Каждая оценка по математике $-$ целое число от $1$ до $5$). Для улучшения оценок он отработал несколько контрольных работ и исправил оценки по ним с «единицы» на «тройку». Докажите, что средняя оценка Незнайки не достигла «четверки».

\problem{6}
В таблице с $4$ строками и k столбцами записаны натуральне числа так, что числа в каждой строке различные и сумма чисел в каждом столбце равна $30$. Найдите максимально возможное количество столбцов в таблице.